\title{Byte Latent Transformer: Patches Scale Better Than Tokens}

\begin{abstract}
We present the Byte Latent Transformer (\textsc{BLT}), a novel architecture for byte-level large language models (LLMs) that achieves performance comparable to traditional tokenization-based models at scale, while offering notable gains in inference speed and robustness. \textsc{BLT} operates by encoding input bytes into adaptive, variable-length patches that become the main units of processing. These patches are determined according to the entropy of upcoming bytes, dedicating greater model attention and computation to regions of higher complexity within the data. We conduct the inaugural flop-controlled scaling analysis of byte-level architectures, training models up to 8B parameters on 4T bytes. The findings affirm that it is viable to scale models directly on byte sequences without depending on a fixed vocabulary. Both the training and inference processes benefit in efficiency, as the model utilizes longer patches when the input exhibits predictability, leading to improvements not only in computation but also in reasoning abilities and generalization to rare or complex cases. Ultimately, given constant inference budgets, \textsc{BLT} achieves superior scaling behavior compared to models with explicit tokenization by enabling simultaneous expansion in both patch and model capacity.
\end{abstract}

\section{Introduction}
\label{section:intro}

We present the Byte Latent Transformer~(\textbf{BLT}), an architecture that dispenses with tokenization and directly leverages raw byte sequences, achieving—for the first time—performance on par with established tokenization-based models at large scale, while also substantially enhancing both computational efficiency and robustness (\S\ref{sec:robustness}).

Current large language models (llms) generally operate in an almost fully end-to-end fashion, with the exception of tokenization. Tokenization, a rule-based pre-processing method, partitions input bytes into a predetermined token inventory, which introduces compression biases and presents several challenges. These include increased sensitivity to specific domains or modalities~\citep{dagangetting}, vulnerability to input perturbations~(\S\ref{sec:robustness}), insufficient orthographic representation~\citep{edman2024cute}, and issues surrounding multilingual fairness~\citep{liang2023xlm,petrov2024language,limisiewicz2024myte}.

Tokenization has historically been crucial because training large language models on raw byte sequences incurs substantial computational expense, primarily because of the increased sequence lengths~\citep{xue2022byt5}. Previous research has attempted to address this challenge by introducing more efficient self-attention mechanisms~\citep{el2020characterbert, clark2022canine} or by leveraging architectures that do not rely on attention~\citep{wang2024mambabyte}~(\S\ref{section:related_work}). Nonetheless, such approaches are predominantly beneficial for \textit{small models}. At large scales, it is the extensive feed-forward layers within Transformers, which are executed over each byte, that account for the bulk of computational cost, overshadowing the contribution of the attention component.

For effective compute allocation, we introduce an adaptive, learnable approach for aggregating bytes into \textit{patches}~(\S\ref{section:patching}), along with a novel model architecture that integrates both byte-level and patch-level representations. In contrast to tokenization, {\textsc BLT} does not impose a predefined patch vocabulary. Instead, it uses lightweight, trainable encoder and decoder components to transform arbitrary byte groupings into corresponding latent patch embeddings. Our experiments demonstrate that this strategy enables a \textit{higher} degree of compute efficiency compared to models relying on traditional tokenization.

In tokenization-based llms, each token is processed with equal computational effort. This approach prioritizes performance at the expense of efficiency, as the tokenization process relies on compression strategies that do not always align with the actual complexity of predicting each token. Our architectural approach is grounded in the principle that models ought to allocate computational resources adaptively, focusing more effort where it is most warranted. For instance, predicting the final letters of most words does not require a large transformer due to the relatively simple and low-entropy nature of these tasks, especially when contrasted with generating the first token of a new sentence, which involves greater uncertainty. In the design of {\textsc BLT} (\S\ref{section:architecture}), this idea manifests as an arrangement of three transformer modules: a pair of compact, byte-level \textit{local models} and a sizable global \textit{latent transformer} (\autoref{fig:arch_high_level}). To inform the formation of byte patches and thus tailor the computational distribution, {\textsc BLT} partitions sequences according to the entropy involved in next-byte prediction, resulting in context-aware byte clusters characterized by more consistent informational content.

We introduce the inaugural study analyzing the scaling properties of byte-level models governed by computational cost, extending up to 8B parameters and 4T bytes of training data. Our results demonstrate that it is feasible to train models end-to-end directly from raw bytes at large scales, bypassing the need for fixed-vocabulary tokenization. In terms of training efficiency, {\textsc BLT} achieves comparable flop-controlled performance\footnote{Model computational cost is assessed as the number of Floating Point OPerations (flops) required.} to Llama 3, while requiring up to 50\% fewer flops during inference~(\S\ref{section:scaling}). 

Moreover, processing raw bytes leads to notable gains in modeling infrequent or atypical data distributions. {\textsc BLT} models show greater robustness to noisy input compared to models that rely on tokenization, and they outperform in tasks demanding deep character-level understanding, such as orthographic competence, phonological analysis, and low-resource machine translation~(\S\ref{sec:robustness}). 

Additionally, {\textsc BLT} enables the joint scaling of both model size and patch size under a constant inference flops constraint. Employing longer patch sizes tends to reduce the computational overhead, thereby permitting an increase in the dimensions of the global latent transformer by virtue of less frequent execution. Our inference-flop-constrained scaling investigations~(\autoref{fig:fixed_inference_scaling}) reveal considerably improved scaling behaviors compared to conventional tokenization-based models.

To conclude, this work presents several key contributions: 1) We propose {\textsc BLT}, a byte latent llm design that adaptively allocates computational resources to enhance flop efficiency; 2) We demonstrate that our approach matches the training flop-controlled performance of Llama 3 at scales up to 8B parameters, and allows for significant trade-offs—achieving up to a 50\% increase in flop efficiency at the cost of only minor decreases in evaluation metrics; 3) The {\textsc BLT} framework enables a new scaling pathway for llms, allowing model parameters to be increased without exceeding a predetermined inference cost; 4) We provide evidence of {\textsc BLT} models’ greater resilience to noisy inputs and their sensitivity to sub-word information in inputs, which standard token-based llms frequently overlook. Our training and inference code for {\textsc BLT} is publicly available at~\url{https://github.com/facebookresearch/blt}.

\section{Patching: From Individual Bytes to Groups of Bytes}
\label{section:patching}

Dividing bytes into \textit{patches} enables {\textsc BLT} to allocate computational resources in a context-sensitive manner. As illustrated in Figure~\ref{fig:patching_types}, there exist several strategies for grouping bytes into patches. More precisely, a patching function $f_p$ takes an input byte sequence $\pmb{x}=\{x_i,|i=1,\ldots n\}$ of length $n$ and produces an output sequence of $m < n$ patches $\pmb{p}=\{p_j|j=1,\ldots,m\}$, assigning each $x_i$ a value from \{0,1\}, where 1 marks the initiation of a new patch. In both token-based and patch-based architectures, the overall computation required to process an input predominantly hinges on the number of iterations the principal Transformer component must perform. For {\textsc BLT}, this figure is directly tied to the total patches the data is split into, as determined by the applied patching function. Therefore, the typical or average length of a patch, referred to as the \textit{patch size}, plays a central role in establishing the compute demands—during both learning and inference—when employing a specific patching approach~(\S\ref{section:flops}). 

In the following sections, we present three distinct patching techniques: fixed-length patching~(\S\ref{section:static-patch}), patching based on whitespace delimiters~(\S\ref{section:white-patch}), and adaptive patching derived from local byte entropy estimates from a compact language model~(\S\ref{section:dyn-patch}). We conclude by examining the process of incremental patching and drawing distinctions between patching and conventional tokenization~(\S\ref{section:bpe-patch}).

\subsection{Strided Patching Every K Bytes}
\label{section:static-patch}

A common approach to organizing bytes is to divide them into fixed-size segments of length $k$, following the method utilized in MegaByte~\citep{yu2023megabyte}. Employing a constant stride facilitates straightforward implementation during both training and inference, enables direct adjustment of the typical patch size, and thereby offers a simple way to regulate computational cost. Nonetheless, this static patching strategy presents notable limitations. Primarily, computation is not adaptively distributed: excessive transformer resources might be spent predicting low-information regions such as whitespace in code, while complex, information-rich regions like mathematical notation may receive insufficient computation. Additionally, this approach causes variability in how similar byte sequences are partitioned, which can result in non-contextual and inconsistent splitting—for example, the segmentation of identical words may differ across occurrences.

\subsection{Space Patching}
\label{section:white-patch}

\citet{slagle2024spacebyte} introduces an enhancement to traditional strided patching by generating new patches immediately after any space-like bytes\footnote{
    Space-like bytes refer to any byte that is neither a Latin letter, a digit, nor a UTF-8 continuation byte; each patch is also required to have at least one non space-like byte.
}, which frequently coincide with natural linguistic segmentations in many languages. The Space patching strategy dedicates a latent transformer computation step (i.e., increases the number of FLOPs) to each word, thereby ensuring consistent patching of words across input sequences and focusing computational resources on challenging predictions that typically arise immediately after spaces. For instance, generating the initial byte for the answer to the query ``Who composed the Magic Flute?\uline{\hspace{.6em}}'' is considerably more complex than generating subsequent bytes after the first letter ``M,'' since the range of potential completions narrows dramatically with that initial character, making it easier to predict the remainder of the answer ``Mozart.'' Nevertheless, the approach of space patching faces difficulties when applied across all languages and domains, and, crucially, it lacks the capacity to dynamically adjust patch sizes. In what follows, we present a novel patching technique inspired by the observation that initial bytes in words are often the hardest to predict, which also offers an intuitive method for patch size control.

\subsection{Entropy Patching: Using Next-Byte Entropies from a Small Byte LM}
\label{section:dyn-patch}

In place of a rule-based method like whitespace segmentation, we adopt a data-driven strategy to detect points of high uncertainty in next-byte prediction. Our approach, termed \textit{entropy patching}, determines patch boundaries through analysis of entropy estimates.

Specifically, we fit a compact auto-regressive language model operating at the byte level on the training data used for {\textsc BLT}, and then evaluate the entropy for the prediction of each subsequent byte according to the language model's distribution $p_{e}$ across the byte vocabulary $\mathcal{V}$:

\begin{equation}
H(x_i) = \sum_{v \in \mathcal{V}} p_{e}(x_i=v|\pmb{x}_{<i}) \log p_{e}(x_i=v|\pmb{x}_{<i})
\end{equation}

We investigate two approaches for detecting patch boundaries using entropies $H(x_i)$. The first method selects points exceeding a predetermined global entropy threshold, as depicted in~\autoref{fig:patching}. The second strategy highlights points where the entropy is significantly higher compared to the immediately preceding value. This latter technique can also be viewed as locating points where the trend of approximately monotonically decreasing entropy within a patch is disrupted.

\begin{align*}
    \text{Global Constraint}\hspace{1cm}H(x_t)& > \theta_g\\
    \text{Approx. Monotonic Constraint}\hspace{1cm}H(x_t)& - H(x_{t-1}) > \theta_r
\end{align*}

Patch boundaries are determined through a simple preprocessing procedure performed at the dataloading stage. This contrasts with the approach of~\citet{nawrot-etal-2023-efficient}, which involves training a classifier to identify patch boundaries based on entropy calculations. In our experimental evaluation~(\S\ref{section:experiments}), we provide a comparison between these two strategies for separating bytes with low and high entropy.

\subsection{The Byte-Pair Encoding (BPE) Tokenizer and Incremental Patching}
\label{section:incremental-patching}
\label{section:bpe-patch}

Numerous contemporary llms, such as the baseline Llama 3 utilized in this work, implement a subword tokenization strategy like bpe~\citep{gage1994new, sennrich-etal-2016-neural}. Throughout this paper, we designate ``tokens'' as byte-based groupings originating from a \textit{finite} vocabulary established before training, in contrast to ``patches,'' which denote adaptively formed sequences lacking a predetermined vocabulary. An essential distinction between these two approaches is that, for tokens, the model is deprived of explicit access to the raw byte-level features.

An important advancement offered by {\textsc BLT} compared to tokenization-based models lies in its novel approach to balancing vocabulary size with computational requirements. Traditional llms encounter a situation where expanding the vocabulary results in tokens representing longer sequences on average, thus reducing the number of model steps. However, this also necessitates a larger output dimension for the model’s final projection layer. This compromise restricts tokenization-based methods from achieving meaningful flexibility with respect to token length and the computational expenses of inference. As an illustration, while Llama 3 manages to boost the mean token length from 3.7 to 4.4 bytes, it does so at the expense of expanding its embedding table fourfold relative to Llama 2.

During generation, {\textsc BLT} must determine at each point in the byte sequence whether it has reached a patch boundary, as this decision dictates if the Latent Transformer should be applied for additional computation. Crucially, this judgment must be made without relying on bytes that have not yet been produced. Therefore, the mechanism for identifying patches cannot presume visibility into upcoming bytes for determining how to segment the sequence. Formally, we define a patching method $f_p$ as possessing the property of incremental patching if it holds that:
$$f_p(\pmb{x}_{<i}) = f_p(\pmb{x})_{<i}$$
Notably, bpe fails to be an incremental patching method, since identical prefixes may be tokenized in divergent ways based on the subsequent sequence, thus violating the property described above\footnote{Insertion of an explicit delimiter token at patch boundaries can adapt bpe to meet the incremental patching requirement; however, this results in a lengthened byte sequence.}.

\section{{\textsc BLT} Architecture}
\label{section:architecture}
{\textsc BLT} integrates a large-scale autoregressive global language model that processes patch-level representations, together with two compact local models that respectively convert byte sequences into patches and reconstruct bytes from these patch representations~(\autoref{fig:arch_high_level}).

\subsection{Latent Global Transformer Model}

\textit{The Latent Global Transformer} refers to an autoregressive transformer model, denoted $\mathcal{G}$, comprising $l_{\mathcal{G}}$ layers. This model transforms an input sequence of latent patch representations, $p_j$, into an output sequence of patch representations, $o_j$. In this work, subscripts $j$ and $i$ consistently refer to patches and bytes, respectively. Employing a block-causal attention mask as introduced by~\citet{dubey2024llama}, the global model limits each patch’s attention scope to the current and preceding patches within the same document. Since this model is responsible for the majority of floating-point operations during both pre-training and inference, selectively applying it enables us to adjust and manage computational expenditure in accordance with the complexity of the input and output regions.

\subsection{Local Encoder}
\label{section:local}

\textit{The Local Encoder Model}, represented as $\mathcal{E}$, is a compact transformer variant with significantly fewer layers ($l_{\mathcal{E}} << l_{\mathcal{G}}$), specifically designed to transform input byte sequences $b_i$ into rich patch-level representations $p_j$ in an efficient manner. Distinct from conventional transformer designs, this model incorporates a cross-attention layer after each transformer layer, enabling it to aggregate byte-level embeddings into patch representations (see~\autoref{fig:crossattn}).

Initially, each input byte $b_i$ is embedded via a lookup in a parameter matrix $\mathbb{R}^{256 \times h_{\mathcal{E}}}$, yielding an embedding $x_i$. Optionally, these embeddings can be enhanced by appending hash-based embeddings as detailed in~(\S\ref{sec:hash-emb}).

Through successive applications of alternating transformer and cross-attention modules~(\S\ref{sec:enc-cross-attn}), these byte-level embeddings are iteratively refined to yield the final patch representations $p_i$, which serve as input to the subsequent global transformer $\mathcal{G}$.

Each transformer block applies a \textit{local block causal} attention mechanism, where each byte has visibility over a sliding context window comprising $w_{\mathcal{E}}$ previous bytes, allowing attention to cross patch borders dynamically (but never document boundaries). The subsections that follow outline specific implementation aspects of the embedding strategy and the design of the cross-attention component.

\subsubsection{Encoder Hash n-gram Embeddings}
\label{sec:ngram-emb}
\label{sec:hash-emb}
An essential element in developing resilient and meaningful representations at each position $i$ involves encoding the context provided by previous bytes. Within {\textsc BLT}, this is realized by representing the byte $b_i$ in isolation as well as in the context of an n-gram of bytes. Specifically, at every step $i$, we begin by forming byte-grams

\begin{align}
    g_{i,n}=\{b_{i-n + 1},\ldots, b_i\}
\end{align}

for each byte position $i$ and $n$ from three to eight.\footnote{
We omit byte-grams of size $n$ or more when $i<n$. }

Next, we present hash $n$-gram embeddings, where every byte $n$-gram is assigned an index in a fixed-size embedding table $E_{n}^{hash}$ through a hash function, for each $n \in \{3,4,5,6,7,8\}$~\citep{bai2010learning}. This embedding is summed with the corresponding byte embedding, normalized, and subsequently used as input for the local encoder. We compute the augmented embedding

\begin{align}
e_i &= x_i + \sum_{n = 3,...,8}E_{n}^{hash}(\text{Hash}(g_{i,n})) \\
\text{where, Hash}(g_{i,n}) &= \text{RollPolyHash}(g_{i,n}) \% |E_{n}^{hash}|
\end{align}

Each $e_i$ is scaled by dividing by the total count of $n$-gram sizes, incremented by one, and we implement RollPolyHash as described in Appendix~\ref{appendix:rollpolyhash}. Section~\ref{section:ablations} presents an ablation study on the impact of varying $n$ and the embedding table dimension for $n$-gram hash embeddings with respect to trends in flop-controlled scaling laws. Beyond hash-based $n$-gram embeddings, we also investigated frequency-based $n$-gram embeddings; further methodological specifics of this approach can be found in Appendix~\ref{appendix:ngrams}.

\subsubsection{Encoder Multi-Headed Cross-Attention}
\label{sec:enc-cross-attn}
Our implementation adopts the input cross-attention approach from the Perceiver architecture~\citep{jaegle2021perceiver}, with the principal modification that the latent vectors here are not fixed but are dynamically constructed as patch-specific representations (see~\autoref{fig:crossattn}). Each patch's representation only interacts with the set of bytes contained within that patch. Specifically, for each patch $p_j$, a query vector is generated by first aggregating the byte representations associated with $p_j$ via a pooling operation, and then applying a linear transformation, $\mathcal{E}_{C} \in \mathbb{R}^{h_{\mathcal{E}} \times (h_{\mathcal{E}} \times U_{\mathcal{E}})}$, where $U_{\mathcal{E}}$ refers to the number of cross-attention heads in the encoder. To put this more formally, let $f_{\text{bytes}}(p_j)$ be the ordered byte sequence for the patch $p_j$. Then, we have:

\begin{align}
P_{0,j} &= \mathcal{E}_{C}(f_\text{bytes}((p_j)), f~ \text{denoting the pooling operation}) \\
P_l &= P_{l-1} + W_o\left(\text{softmax}\left(\frac{QK^T}{\sqrt{d_k}}\right)V\right) \\
\text{where } Q_j &= W_q(P_{l-1,j}), K_i = W_k(h_{l-1,i}), V_i = W_v(h_{l-1,i}) \\
h_l &= \text{Encoder-Transformer-Layer}_l(h_{l-1})
\end{align}

Here, $P \in \mathbb{R}^{n_p \times h_{\mathcal{G}}}$ denotes a set of $n_p$ patch representations that are input into the global module. These patch representations are formed by aggregating the byte embeddings $e_i$ that make up each respective patch $p_j$. The matrices $W_q$, $W_k$, $W_v$, and $W_o$ serve as the learnable linear transformations for queries, keys, values, and outputs, with the key and value vectors produced by projecting the previous layer’s byte representations $h_i$ (or $e_i$ in the initial layer). To ensure localized attention, a specialized masking scheme is applied: each query $Q_j$ is constrained to attend only to those keys and values associated with the bytes inside its corresponding patch $j$. Given that patch representations tend to have a higher dimensionality $h_{\mathcal{G}}$ compared to $h_{\mathcal{E}}$, and since attention utilizes multiple heads across $Q, K, V$, we represent $P_l$ as a collection of multi-head vectors each of size $h_{\mathcal{E}}$ for cross-attention, and subsequently concatenate these head outputs to recover the full $h_{\mathcal{G}}$ dimension. The module further applies pre-LayerNorm normalization separately to queries, keys, and values, and explicitly omits positional encodings in this cross-attention operation. As a final architectural detail, a residual connection encloses the entire cross-attention component.

\subsection{Local Decoder} 

Analogous to the structure of the local encoder, the local decoder $\mathcal{D}$ utilizes a transformer-based architecture characterized by its lightweight nature, comprising $l_{\mathcal{D}} << l_{\mathcal{G}}$ layers. This decoder takes the sequence of global patch representations $o_j$ and reconstructs them into raw bytes $y_i$. It generates each byte in the sequence conditioned on previously generated bytes, leveraging the hidden states output by the local encoder corresponding to the byte sequence. The decoding process consists of $l_{\mathcal{D}}$ repeated blocks, each containing a cross-attention mechanism followed by a transformer layer. Within each block, cross-attention is applied first to convert patch-level representations into byte-level representations, which are then further refined by the subsequent local transformer layer that processes the resulting byte sequence.

\subsubsection{Decoder Multi-headed Cross-Attention} 

Within the decoder’s cross-attention mechanism, the query and key/value assignments are reversed: here, the queries are provided by the byte-representations, while the keys and values are derived from the patch representations. The starting point for the byte-representations in the cross-attention process is established using the byte embeddings output from the encoder’s final layer, i.e., $h_{l_{\mathcal{E}}}$. For each layer $l$, the updated byte-representations $d_{l,i}$ are determined according to:

\begin{align}
D_0 &= h_{l_{\mathcal{E}}} \\
B_l &= D_{l-1} + W_o\left(\text{softmax}\left(\frac{QK^T}{\sqrt{d_k}}\right)V\right), \\
  &\text{with } Q_i = W_q(d_{l-1,i}), K_i = W_k(\mathcal{D}_C(o_j)), V_i = W_v(\mathcal{D}_C(o_j)) \\
  D_l &= \text{Decoder-Transformer-layer}_l(B_l)
\end{align}

Here, $W_k$ and $W_v$ denote the projection matrices for keys and values, which are applied to a linearly transformed and split operation $\mathcal{D}_C$ on the final patch representations $o_j$ produced by the global model. $W_q$ serves as the query projection matrix, applied to the byte-level representations $d_{l-1}$ from the preceding decoder transformer layer (or $h_{l_{\mathcal{E}}}$ when considering the first layer), while $W_o$ projects the resulting outputs. Consequently, the matrix $B$ takes the form $B \in \mathbb{R}^{h_{\mathcal{D}} \times n_b}$, where $n_b$ corresponds to the total number of output bytes. The updated decoder representations $D_l$ are obtained by passing $B$, which is the output from the cross-attention module, through a decoder transformer layer. Consistent with the approach taken in the local encoder cross-attention, the mechanism employs multiple attention heads, utilizes pre LayerNorms, omits positional embeddings, and includes a residual connection that bypasses the cross-attention component.

\section{Experimental Setup}
\label{section:experiments}
We conduct rigorously controlled experiments to directly assess the performance of {\textsc BLT} in comparison with tokenization-based models, ensuring that {\textsc BLT} does not benefit from any potential increases in sequence context length.

\subsection{Pre-training Datasets} In this study, all model variants are initialized with pre-training on two principal datasets: 1) The Llama 2 dataset~\citep{touvron2023llama}, containing 2 trillion tokens sourced from an array of publicly available corpora, which have undergone extensive cleaning and filtering to enhance data quality; and 2) {\textsc BLT}-1T: a newly constructed dataset comprising 1 trillion tokens extracted from diverse public sources, and featuring a subset of pre-training material made available through Datacomp-LM~\citep{li2024datacomp}. The Llama 2 dataset is employed primarily in scaling law analyses, guided by the recommendations in~\cite{dubey2024llama}, to inform optimal architectural configurations for {\textsc BLT}. In contrast, {\textsc BLT}-1T serves as the foundation for a full pre-training cycle, facilitating head-to-head downstream performance evaluations with Llama 3. Neither dataset contains information from Meta platforms or services. For tokenization baselines, experiments utilize the Llama 3 tokenizer, which features a vocabulary of 128K tokens and demonstrates superior baseline results compared to the Llama 2 tokenizer according to our empirical assessments.

\subsection{Entropy Model} For our experiments, the entropy model serves as a byte-level language model, utilizing the same training data as the complete {\textsc BLT} model. Unless specified otherwise, the architecture comprises a transformer equipped with 100M parameters across 14 layers, a hidden size of 512, and employs sliding window attention spanning 512 bytes. All other hyperparameters mirror those used in both our local and global transformer setups. We further explored a variety of model configurations—altering size, receptive field, and architectural choices—as outlined in \autoref{section:ablations}. Notably, when configured with a sufficiently limited receptive field, the entropy model is compact enough to be efficiently represented as a lookup table.

\subsection{Entropy Threshold and Equalizing Context Length} For models employing entropy-based patching, we determine a patching threshold designed to yield a specified mean \textit{patch size} over the mixture of pretraining datasets. In {\textsc BLT}, as opposed to standard tokenization schemes, the \textit{patch size} is flexibly defined and can be set arbitrarily, which substantially impacts the model's effective context window. To ensure comparability and prevent configurations with larger patches from obtaining undue benefit, we fix the expected total byte count per batch across experiments. Consequently, sequence lengths are shortened for models utilizing larger patches. Specifically, we employ an 8k byte context size for Llama 2 data, while the context is increased to an average of 16k bytes for the {\textsc BLT}-1T dataset, with the overall average batch size held constant at 16M bytes.

Although the average batch size remains unchanged, dynamic patching strategies produce varying proportions of bytes per patch when assembling data batches. To enhance efficiency, our approach to {\textsc BLT} training organizes batches so that each one contains an equal number of patches, thereby eliminating the need for padding within the computationally intensive latent transformer. For training purposes, byte sequences are padded or truncated to lengths of 12k and 24k bytes, corresponding to the Llama 2 and {\textsc BLT}-1T datasets, respectively, to prevent memory surges caused by sequences with exceptionally large patches.

\subsection{Entropy Model Context}
\label{section:entropy-context}
Through experimentation, we observe that applying entropy patching results in increasingly large patches, particularly in structured tasks such as multiple choice questions (refer to the MMLU patching example in \autoref{fig:mmlu_patching}), where repetition is frequent. This phenomenon arises from the reduced entropy present in repetitive segments within the entropy model context. Consequently, for the large-scale implementation of {\textsc BLT}-Entropy utilizing a patch size of 4.5, we mitigate "entropy drift"—which stems from varying context lengths—by periodically resetting the entropy context with line breaks and employing an approximate monotonicity constraint. This adjustment solely modifies the entropy calculation approach; our methodology for determining the entropy threshold remains unchanged.

\subsection{FLOPs Estimation} 
\label{section:flops}

Our methodology for calculating transformer flops is based primarily on the formulations presented in Chinchilla~\citep{hoffmann2022training}, incorporating contributions from the feed-forward modules, qkvo projections in the self-attention mechanisms, as well as the computations involved in attention and output projections. However, one key distinction in our approach is the assumption that the input embedding is realized via an optimized lookup operation rather than dense matrix multiplication, resulting in zero computational cost for this layer. Consistent with prior literature, we also approximate the computational requirements for the backward pass as being twice that of the forward computation.

To determine the number of flops \textit{per byte} in {\textsc BLT} models, we aggregate the computational costs from the local encoder transformer, global latent transformer, and local decoder transformer, as well as the cross-attention mechanisms found in both the encoder and decoder:

\begin{align}
    \text{FL}_{\text{{\textsc BLT}}} &= \text{Transf. FL}(h_{\mathcal{G}}, l_{\mathcal{G}}, m=n_{ctx}/n_p, V=0) / n_p \\
    &+ \text{Transf. FL}(h_{\mathcal{E}}, l_{\mathcal{E}}, m=w_{\mathcal{E}}, V=0) \\
    &+ \text{Transf. FL}(h_{\mathcal{D}}, l_{\mathcal{D}}, m=w_{\mathcal{D}}, V=256) \\
    &+ \text{Cross Attn. FL}(h_{\mathcal{E}}, l_{\mathcal{E}}, m=n_p, r=n_p/k) \times k/n_p \\
    &+ \text{Cross Attn. FL}(h_{\mathcal{D}}, l_{\mathcal{D}}, m=k, r=k/n_p)
\end{align}

Here, $n_{ctx}$ denotes the length of the input sequence measured in bytes, $n_p$ indicates the size of each patch, $r$ represents the proportion of query vectors to key/value vectors, and $k$ quantifies the relationship between the patch dimension and the byte dimension—specifically, it refers to the number of local model segments concatenated to construct the complete model representation (in \autoref{fig:crossattn}, $k=2$). $V$ denotes the size of the output vocabulary, utilized exclusively within the local decoder. The value of $m$ signifies the context length in the attention mechanism and varies depending on whether the module processes byte-level or patch-level sequences. We adapt the attention-related flop calculations based on the role and structure of each component. The specific formulas used for calculating Transformer-FLOPs and Cross-Attention FLOPs are outlined in Appendix~\ref{section:flops-appendix}.

\subsection{Bits-Per-Byte Estimation} Perplexity is meaningful solely when the tokenizer is held constant, as it quantifies the predictive uncertainty of individual tokens. In the evaluation of models operating at the byte versus token level, prior research~\citep{xue2022byt5, yu2023megabyte, wang2024mambabyte} has adopted Bits-Per-Byte (BPB) as the evaluation metric. BPB serves as a version of perplexity that does not depend on any specific tokenizer. Concretely:

\begin{align}
    \text{BPB}(x)&= \frac{\mathcal{L}_{CE}(\pmb{x})}{\ln(2)\cdot n_{\text{bytes}}}
\end{align}

where the uncertainty over the data $\pmb{x}$ as measured by the sum of the cross-entropy loss is normalized by the total number of bytes in $\pmb{x}$ and a constant. 

\subsection{Transformer Architecture Hyperparameters} 

For each transformer block within {\textsc BLT}, encompassing both local and global configurations, our setup closely mirrors that of Llama~3~\citep{dubey2024llama}. Specifically, we implement the SwiGLU activation function~\citep{swiglu} within all feed-forward networks, utilize rotary positional embeddings (RoPE)~\citep{RoPe} parameterized with $\theta=500000$~\citep{xiong2024effective} in the self-attention modules alone, and employ RMSNorm \citep{RMSNorm} for normalizing layers. For self-attention blocks utilizing standard attention masks, namely \textit{block causal} and \textit{fixed-window block causal} masks, we adopt Flash Attention~\citep{dao2022flashattention}. When applying fixed-width attention masks, we set the window size at 512. The cross-attention layers, which require dynamically determined patch-based masks, are equipped with Flex Attention\footnote{\url{https://pytorch.org/blog/flexattention}}. This choice leverages fused kernel implementations, resulting in substantial training speed-ups.

\subsection{{\textsc BLT}-Specific Hyperparameters}
To evaluate the performance of {\textsc BLT} models, our experimental setup explores both scaling patterns and performance on downstream tasks. We examine models of varying capacities: 400M, 1B, 2B, 4B, and 8B parameters for this analysis. Details regarding their architectural hyperparameters can be found in Appendix~Table~\ref{tab:arch}.

Initialization of queries in the initial cross-attention layer within the local encoder is performed via max-pooling. Each {\textsc BLT} model employs a set of $500,000$ hashes with a single hash function and uses n-gram sizes between 3 and 8. A consistent learning rate of $4e-4$ is used across all models, a decision supported by a hyperparameter sweep over $1e-3$ to $1e-4$ at both the 400M and 1B parameter scales, which indicated that this value is ideal for both token and {\textsc BLT} model variants. For experiments following the scaling trends based on Llama-2 data, we adopt the batch-size recommendations from~\cite{dubey2024llama} or the corresponding byte-based equivalent.

Optimization is carried out with the AdamW optimizer~\citep{adamw}, setting $\beta_1$ to 0.9, $\beta_2$ to 0.95, and $\epsilon=10^{-8}$. The training procedure involves 2000 steps of linear warm-up, followed by a cosine learning rate decay to zero. Additionally, a weight decay factor of 0.1 is enforced, and global gradient norms are clipped at 1.0.

\section{Scaling Trends}
\label{section:scaling}

We offer a comprehensive analysis of the scaling behaviors exhibited by byte-level models, providing valuable insights for future scaling of {\textsc BLT} models. Our investigation into scaling extends beyond the scope of earlier studies on byte-level approaches by: (a) contrasting performance across the compute-optimal training paradigm, (b) training substantial 8B-parameter models with significant data volumes (up to 1T tokens/4T bytes) and assessing them on downstream benchmarks, and (c) evaluating scaling dynamics under scenarios where inference cost is held constant. A subsequent section will delve into particular benefits unique to modeling byte sequences.

\subsection{Parameter Matched Compute Optimal Scaling Trends}
\label{section:parameter-matched-scaling}
Leveraging the Llama 2 dataset, we conduct training of several \textit{compute-optimal} bpe and {\textsc BLT} models at four distinct model capacities, spanning from 1B to 8B parameters. Training FLOPs are then charted against language modeling accuracy on a selected portion of the data mixture. The bpe models utilize the empirically optimal parameter-to-data ratio established by Llama~3~\citep{dubey2024llama}. This \textit{compute-optimal} regime is intended, in theory, to deliver maximal performance on the training set for a given computational expenditure~\citep{hoffmann2022training}, thereby serving as a strong comparison point for our own models. For each bpe configuration, we also develop an identically parameterized {\textsc BLT} model on the same dataset, employing a Latent Transformer whose architecture precisely mirrors its bpe counterpart.

As shown in~\autoref{fig:scaling-compute-opt} (right), {\textsc BLT} models consistently achieve results that are on par with or superior to those of bpe models, with this pattern persisting across various scales of model size and computational cost. According to our understanding, {\textsc BLT} is the first Transformer architecture at the byte level to demonstrate scaling properties comparable to BPE-based systems under compute-optimal conditions. This observation supports our hypothesis that the parameter-to-compute ratio deemed optimal for bpe is also applicable to {\textsc BLT}, or at the very least, the ratios are closely aligned.

Both enhancements in architecture and the use of dynamic patching play essential roles in aligning with bpe scaling behaviors. In ~\autoref{fig:scaling-compute-opt} (left), space-patching-based approaches are evaluated alongside Llama 3. For this purpose, SpaceByte \citep{slagle2024spacebyte} is approximated via {\textsc BLT} space-patching, excluding both n-gram embeddings and cross-attention mechanisms. While SpaceByte demonstrates gains over Megabyte, it still does not reach the performance levels of Llama 3. The panel on the right in \autoref{fig:scaling-compute-opt} showcases the combined impact of both architectural refinements and dynamic patching; at scale, {\textsc BLT} models achieve results that rival those of leading tokenizer-based systems such as Llama 3.

Furthermore, our analysis reveals that the selection of a tokenizer significantly affects the outcomes of tokenizer-based models. Specifically, models trained with the Llama-3 tokenizer show superior performance compared to those trained with the Llama-2 tokenizer, given the same training dataset.

In summary, our {\textsc BLT} model occupies a middle ground between Llama 2 and Llama 3 when evaluating performance at much larger patch sizes. The average token sizes produced by the bpe tokenizers in Llama 2 and 3 are 3.7 and 4.4 bytes, respectively. On the other hand, {\textsc BLT} demonstrates the ability to reach comparable scaling dynamics even when average patch sizes reach 6 or 8 bytes. Given that inference flops decrease as average patch size increases, adopting an 8-byte patch size can result in roughly a 50\% reduction in inference compute. Moreover, it appears that scaling both the model and data size further improves the effectiveness of larger patch sizes. Notably, when employing an 8-byte patch size, {\textsc BLT} initially underperforms relative to bpe Llama 2 at the 1B parameter level but ultimately surpasses bpe Llama 2 at the 7B parameter level. This pattern indicates that performance advantages at high patch sizes may become even more pronounced with continued scaling, raising the possibility that even greater patch sizes could be advantageous with further increases in model capacity and training computation.

\subsection{Beyond Compute Optimal Task Evaluations}
\label{section:evals}
To further explore scaling dynamics, we extend training of an 8B {\textsc BLT} model beyond the conventional compute-optimal regime, utilizing the {\textsc BLT}-1T dataset—a larger, higher-quality dataset—to observe its behavior on a variety of established classification and generation benchmarks. For this analysis, we focus on a selection of tasks assessing commonsense reasoning, general world knowledge, and programming proficiency:
\paragraph{Classification tasks} considered are ARC-Easy (0-shot)~\citep{Eval_ARC}, Arc-Challenge (0-shot)~\citep{Eval_ARC}, HellaSwag (0-shot)~\citep{Eval_hellaswag}, PIQA (0-shot)~\citep{Eval_piqa}, and MMLU (5-shot)~\citep{hendrycks2020measuring}. We utilize a prompt-based approach, evaluating the predicted likelihood of response options, and present the mean accuracy across these datasets.
\paragraph{Coding related generation tasks:} We provide pass@1 results for MBPP (3-shot)~\citep{Eval_MBPP} and HumanEval (0-shot)~\citep{Eval_HumanEval} as a measure of the model’s Python code synthesis capabilities.

In \autoref{tab:evals}, we present a comparative analysis of three models trained on the {\textsc BLT}-1T dataset: a model built using the bpe Llama 3 tokenizer,\footnote{We utilize the Llama 3 tokenizer, featuring a 128k vocabulary, as it surpasses the Llama 2’s 32k vocabulary in performance.} along with two distinct {\textsc BLT} model variants. The first variant applies a space-patching method ({\textsc BLT}-Space), while the second implements an entropy-driven patching strategy ({\textsc BLT}-Entropy), incorporating an approximate monotonicity constraint and context resets for the entropy-based model triggered by new lines (per the procedure described in~\autoref{section:entropy-context}). Each model is allocated a comparable number of flops during training. For the {\textsc BLT}-Entropy model, we further introduce an inference-time modification, adjusting the entropy threshold from 0.6 down to 0.1, which we observe leads to enhanced task outcomes, albeit requiring a greater number of inference steps.

The {\textsc BLT}-Entropy model achieves superior results compared to the Llama 3 model in 4 of the 7 evaluated tasks, despite both being trained on an equal amount of data. This performance gain likely stems from two primary factors: (1) a more efficient utilization of training resources afforded by dynamic patching, and (2) the explicit representation of information at the byte level rather than the token level.

Conversely, {\textsc BLT}-Space generally lags behind the Llama 3 tokenizer across most tasks, except for one, yet it manages to considerably lower inference flops due to its higher average patch size of 6 bytes. By contrast, both bpe and entropy-patching based approaches exhibit similar average patch sizes, each around 4.5 bytes on the mixed training data. Given identical training budgets, the model employing larger patch sizes processes 30\% more data than the other two, which could potentially cause {\textsc BLT} to deviate further from the compute-optimal threshold.

\subsection{Patches Scale Better Than Tokens} Utilizing {\textsc BLT} models allows for the concurrent scaling of both model and patch sizes, all without increasing the training or inference computational costs or changing the volume of training data. The ability to arbitrarily enlarge patch size is distinctive to patch-oriented approaches and liberates them from the efficiency constraints typically encountered by models based on a fixed-vocabulary tokenization scheme, as elaborated in Section~\ref{section:bpe-patch}. Extending patch length leads to computational savings, permitting those resources to be redirected toward expanding the capacity of the global latent transformer since its application frequency decreases.

We perform a fixed inference scaling analysis to investigate whether increasing model size—while decreasing the number of steps over larger patches—yields improved performance compared to smaller models utilizing more steps. Utilizing the Llama 2 tokenizer, we begin with models containing 400m and 3.6B parameters, and then identify flop-equivalent counterparts built with the Llama 3 tokenizer as well as {\textsc BLT}-Entropy models configured for average patch sizes of 6 and 8 bytes across the training data mixture (refer to~\autoref{tab:fixed-inf-params} for detailed model configurations). For the patch size 8 variants, we employ 3 encoder layers in lieu of 1. All models are trained across a range of training flop budgets.

\autoref{fig:fixed_inference_scaling} demonstrates that {\textsc BLT} models consistently exhibit superior scaling behavior compared to tokenization-based methods across both inference FLOP classes. Initially, BPE models demonstrate better performance when training resources are limited, but are soon overtaken by {\textsc BLT} as training computation slightly exceeds the compute-optimal threshold. Consequently, it can be advantageous to allocate additional resources to one-off pretraining in order to obtain models with superior accuracy at a predetermined inference cost. A notable instance of this phenomenon is seen in 8B-parameter models, such as Llama 3.1, which are pretrained on datasets two orders of magnitude larger than the compute-optimal amount for their size.

The threshold at which {\textsc BLT} surpasses token-based models has moved incrementally nearer to the compute-optimal point with the transition to models in the higher flop class, decreasing from approximately 3x to 2.5x the optimal compute budget. Additionally, the model utilizing a patch size of 8 demonstrates a sharper scaling trajectory within the larger flop class, surpassing alternative models at an earlier stage. As elaborated in Section~\ref{section:parameter-matched-scaling}, increasing patch sizes tends to yield performance outcomes more comparable to those achieved by BPE models as the model scale increases. This trend can be partly explained by the diminishing proportion of total flops allocated to the byte-level Encoder and Decoder components, which exhibit slower scaling behavior in contrast to the Latent Transformer. For instance, when expanding total parameters by a factor of 20, from 400M to 8B, the local parameters of {\textsc BLT} only increase by about twofold. This distinction is significant, given that the use of larger patch sizes primarily affects the computation required by the patch Latent Transformer, leaving the byte-level modules largely unaffected. As a result, {\textsc BLT}-Entropy with patch size 8 experienced an increase from 1.6x to 1.7x the model size of Llama 2 when the model scale was expanded.

To summarize, our investigation into patch-length scaling reveals that the {\textsc BLT} patch-based architecture benefits from enhanced scaling behavior when both patch size and model capacity are increased together. These favorable scaling patterns are maintained and may even become more pronounced at larger model scales.

\section{Byte Modeling Improves Robustness} We further assess the robustness of {\textsc BLT} in comparison to token-based models that do not utilize explicit byte-level inputs, and introduce a methodology for converting pretrained token-based models to operate at the byte level.
\label{sec:robustness}

\subsection{Character-Level Tasks}

Initial efforts to develop byte-level models were primarily driven by their increased resilience to noise introduced at the byte level in the input, as well as their innate understanding of token substructure—capabilities that often challenge models relying on tokenizers. To systematically assess these advantages, we conduct supplementary evaluations using benchmarks specifically designed to test robustness against input noise and to assess sensitivity to character-level details across languages, encompassing not only English, but also a range of multilingual tasks involving digits and phonemes. The outcomes of these analyses are reported in Table~\ref{tab:char_tasks}.

\paragraph{Noisy Data} To assess the resilience of tokenizer-based models in comparison to {\textsc BLT}, we generate altered versions of the benchmark classification datasets outlined in Section~\ref{section:evals}. Five different character-level noise methods are employed to modify the input text: (a) \textit{AntSpeak}: The entire text is rendered in uppercase with spaces inserted between each character. (b) \textit{Drop}: Randomly deletes 10\% of the text's characters. (c) \textit{RandomCase}: Randomly assigns uppercase to half the characters and lowercase to the other half. (d) \textit{Repeat}: Selects 20\% of the characters and repeats them up to four times each. (e) \textit{UpperCase}: Changes every character in the text to uppercase. For each evaluation, the noise is applied either to the prompt, the completion, or both, creating distinct test conditions; scores are then averaged across these. The results for the noised HellaSwag~\citep{Eval_hellaswag} tasks, as shown in Table~\ref{tab:char_tasks}, demonstrate that {\textsc BLT} consistently surpasses tokenizer-based models in terms of robustness—achieving an average improvement of 8 points relative to the model trained on identical data—and even outperforms Llama 3.1, which was trained using a substantially larger dataset.

\paragraph{Phonology - Grapheme-to-Phoneme (G2P)} We evaluate the effectiveness of {\textsc BLT} in converting a sequence of graphemes (written characters forming a word) into the corresponding sequence of phonemes (symbols denoting pronunciation). Table \ref{tab:char_tasks} displays our findings for the G2P task performed in a 5-shot scenario, utilizing the Phonology Bench~\citep{suvarna2024phonologybench}. Our analysis shows that {\textsc BLT} surpasses the baseline performance of the Llama 3 1T tokenizer-based model on this task.

\paragraph{CUTE}
To measure character-level comprehension, we evaluate {\textsc BLT} using the CUTE benchmark~\citep{edman2024cute}, which consists of a range of tasks grouped into three main areas: compositional understanding, recognition of orthographic similarity, and sequence manipulation skills. This suite is notably difficult for many tokenizer-based models, as they tend to recognize the spelling of their tokens yet are generally inadequate at leveraging this knowledge for text manipulation. As indicated in Table~\ref{tab:char_tasks}, {\textsc BLT}-Entropy surpasses both BPE Llama 3 models by over 25 points on the benchmark. Our approach shows particularly strong competence on character manipulation subtasks, attaining a 99.9\% score on both spelling-related tasks. These sizable gains—achieved even though {\textsc BLT} was exposed to 16x less training data than Llama 3.1—suggest that BPE models have difficulty internalizing character-level patterns. Figure \ref{fig:cute_examples} presents examples highlighting instances where the Llama 3 tokenizer model falters but {\textsc BLT} succeeds. Only on the word deletion and insertion tasks do BPE models outperform; while manipulating words may not be straightforward for a byte-level model, the performance disparity is minimal, implying that learning to represent words from characters is possibly less challenging than the reverse. We retain the same experimental settings across all evaluations, utilizing the original Huggingface prompts. It is worth noting that BPE models could potentially see improvement from refined prompt design.

\paragraph{Low Resource Machine Translation} We assess the performance of {\textsc BLT} on translation tasks involving both directions—into and out of English—across six widely represented language families and twenty-one low-resource languages employing diverse scripts, as defined in the FLORES-101 benchmark~\citep{goyal2022flores}. SentencePiece BLEU scores summarizing these outcomes are presented in Table~\ref{tab:flores}. The evaluation reveals that {\textsc BLT} surpasses a baseline model using the Llama 3 tokenizer, offering a 2-point gain for translation into English and a 0.5-point improvement for translation out of English. For high-frequency language pairs, the performance of {\textsc BLT} is on par with, or marginally better than, Llama 3. Significantly, for many pairs among low-resource languages, {\textsc BLT} consistently outperforms Llama 3, highlighting byte modeling's enhanced capability to generalize effectively to underrepresented or rare byte sequences.

\subsection{Training {\textsc BLT} from Llama 3}
\label{sec:distillation}
We investigate a methodology wherein {\textsc BLT} models benefit from pre-existing tokenizer-based models by initializing their global transformer parameters with those extracted from a pre-trained Llama 3.1. This strategy accelerates training and enhances convergence. In this framework, the global transformer’s weights are updated at a rate one-tenth that of the local encoder and decoder models, while maintaining the optimal number of Llama 3 training steps. Table~\ref{tab:distill} presents a head-to-head comparison with a baseline {\textsc BLT}. Results show that {\textsc BLT} initialized from Llama 3.1 achieves substantial performance improvements, surpassing both the original Llama 3 and standard {\textsc BLT}—all trained under equivalent computational budgets. Notably, even against our {\textsc BLT}-Entropy variant (detailed in Table~\ref{tab:evals}), which was exposed to a much larger 1T-token training set, the Llama 3.1-based {\textsc BLT} attains higher scores on MMLU, underscoring this transfer procedure's promise for greatly lowering the training flops required.

This approach can be interpreted as converting tokenizer-based architectures into ones that operate without tokenizers, thereby transforming a pre-trained LLaMA 3.1 into a {\textsc BLT} model. For a thorough assessment, we report results for the original LLaMA 3.1 model—trained with 15T tokens—in Table \ref{tab:distill}, and directly compare its performance with the {\textsc BLT} variant built from LLaMA 3. Our findings indicate that while there is only a modest decrease in scores on MMLU and HumanEval, performance drops more notably on several other benchmarks. These results indicate that additional research is necessary to more effectively utilize the pre-trained base, with particular attention to refining data mixtures and tuning hyperparameters for optimal results.

\section{Ablations and Discussion}
\label{section:ablations}
Here, we present ablation studies that provide support for key architectural decisions made in {\textsc BLT}, alongside our selected patching strategy and hyper-parameters for the {\textsc BLT} model with 8B parameters, trained using the {\textsc BLT}-1T dataset.

\paragraph{Entropy Model Hyper-parameters} To evaluate the impact of the entropy model’s scale and context window length on downstream scaling behavior, we experiment with byte-level entropy transformer architectures spanning a range from 1 million up to 100 million parameters and context window lengths of 64 through 512. Scaling curves of bits per byte (bpb) versus training FLOPs for models with $400m$ and $1b$ parameters trained on the Llama-2 dataset are shown in \autoref{fig:entropy_ablation}. Our results indicate that increased model size and longer context windows improve scaling performance, though the gains plateau for model sizes exceeding 50 million parameters.

\paragraph{Types of Patching}

We conduct ablation studies on the four distinct patching methods outlined in Section~\ref{section:patching}: (1) Strided Patching employing strides of 4 and 6, (2) Whitespace-based Patching, (3) BPE Tokenizer patching utilizing the Llama 3 tokenizer, and (4) Entropy-based patching leveraging a compact byte-level language model.

Although dynamic patching lowers the effective sequence length, we regulate sequence lengths to ensure a consistent context window across all patching methods. Each model is thus exposed to an equivalent number of bytes per sequence throughout both training and inference on average, thereby eliminating confounding effects related to modeling broader contexts. Figure~\ref{fig:scaling-compute-opt} presents the outcomes of these comparative experiments. Every alternative patching method surpasses static patching, with space patching demonstrating performance nearly matching that of dynamic entropy based patching.

\autoref{tab:evals_ablation} provides a comparison of {\textsc BLT} models evaluated on benchmark tasks, focusing on models utilizing tokenizers, space patching, and entropy-based patching. These models were all trained on the Llama 2 dataset for the optimal number of training steps as outlined by \citet{dubey2024llama}. While space patching offers a more straightforward approach by avoiding the computational overhead of an entropy model during training, our results indicate that the advantages of entropy-based patching observed in scaling experiments~(Section~\ref{section:scaling}) also extend to these benchmark evaluations.\footnote{For space patching, we show results from earlier experiments that did not include cross-attention; however, comparable patterns were found even when cross-attention was implemented.}

\paragraph{Cross-Attention} 
In~\autoref{tab:crossattn_ablations_1b}, we conduct an ablation study on the placement of cross-attention modules within both the encoder and decoder of {\textsc BLT}. For the encoder, we explore three different methods to initialize the queries for cross-attention: 1) a shared learned embedding applied uniformly across all global states, 2) embeddings generated by hashing the individual patch bytes, and 3) representations obtained by pooling the encoder’s hidden states corresponding to the patch bytes at the target encoder layer.

Our results indicate that incorporating cross-attention within the \textit{decoder} yields the greatest benefit. When implemented in the encoder, cross-attention shows a minor positive effect, which is observed exclusively when queries are initialized using pooling. Furthermore, cross-attention proves especially advantageous on Common-Crawl data, with the effect being most pronounced for larger patch sizes.

\paragraph{n-gram Hash Embeddings} 

We experiment with n-gram hash embedding vocabulary sizes of 0, 100K, 200K, and 400K, with corresponding results summarized in Table~\ref{tab:ngram_ablations_1b}. Our analysis demonstrates that incorporating hash embeddings yields improvements across all evaluated domains, with the largest impact seen on Wikipedia and Github (yielding a 0.04 bpb benefit versus only 0.01 bpb following 15k training steps at the 8B parameter scale). Adjusting the number of hashes from 500K to 300K at the 8B scale influences the outcome by only approximately 0.001 bpb after 15k steps. These findings suggest that hashing plays a crucial role in aligning the performance of {\textsc BLT} with that of models based on tokenization, yet performance gains plateau beyond 300K hashes. Moreover, improvements from hash embeddings and cross-attention seem to be largely additive, as each contributes to better results on distinct datasets.

\paragraph{Local Model Hyperparameters} Table \ref{tab:local_layers} presents an ablation study over different configurations for the number of layers in both the local encoder and decoder. Our results indicate that when utilizing hash n-gram embeddings, {\textsc BLT} achieves strong performance with a highly minimal encoder—specifically, using only a single layer—while benefitting from a comparatively deeper decoder.

\section{Related Work}
\label{section:related_work}

\textbf{Character-Level RNNs:} The task of Character Language Modeling has long stood as a cornerstone in neural network research~\citep{sutskever2011generating,mikolov2012subword,graves2013generating}, largely due to its capacity for naturally handling out-of-vocabulary words, thus avoiding the need for traditional back-off strategies. In the approach presented by \cite{kim2016character}, models exclusively process character sequences on the input side via convolutional and highway architectures, which then connect to LSTM-based RNNs. Their models demonstrate parity with contemporary RNN-based state-of-the-art language models on English and yield superior results on languages with complex morphology—highlighting a major strength of character-level LLMs. Expanding on this, \cite{kenter2018byte} investigated byte-level LSTM architectures for machine comprehension, obtaining results that exceeded those of word-level counterparts on languages with intricate morphology such as Turkish and Russian. In a related effort, \cite{zhang2015character} leveraged character-level convolutional neural networks for text classification tasks and found them to outperform word-based models in specific applications. Hierarchical LSTM frameworks by \citet{chung2019hierarchical} use boundary detectors at different levels to uncover implicit text hierarchies, contributing further improvements in character-level language modeling performance. Distinctively, ByteNet~\citep{kalchbrenner2016neural} employs convolutional networks over characters for machine translation, diverging from attention-based mechanisms.

\textbf{Character-Level Transformers:} The introduction of transformer architectures utilizing attention mechanisms~\citep{vaswani2017attention}, in conjunction with subword tokenization techniques~\citep{sennrich-etal-2016-neural}, has greatly advanced the capabilities of neural networks on tasks such as language modeling and standard benchmarks. Despite this progress, segmenting text into word or sub-word units imposes an inherent inductive bias on the abstraction level at which models operate. Aiming to bridge transformer advancements with earlier promising character-level results, \citet{al2019character} implemented substantially deep transformer networks and leveraged auxiliary loss functions to train transformer-based models that surpassed character-level LSTM language models. Nonetheless, a notable performance discrepancy remained in comparison to word-level LLMs. Similarly, GPT-2~\citep{radfordgpt2} reported that byte-level language models did not perform as well as word-level models when evaluated on large datasets, such as the 1 billion word benchmark.

\cite{Choe2019BridgingTG} showed that transformer-based byte-level LLMs can surpass subword-level models of similar size, though this comes at a substantial increase in computational requirements and training time. In a related effort, \cite{el2020characterbert} introduced CharFormer, a BERT variant that constructs word-level representations through convolutions on character embeddings. Although this method leads to enhanced performance within the medical domain, it similarly demands significantly greater computational resources. In another study, \cite{clark2022canine} presented CANINE, an encoder-only model with 150M parameters that processes raw character sequences. CANINE incorporates a deep transformer architecture resembling the global module described in our work, paired with local transformers and strided convolutions to reduce input character sequence length. This design allows CANINE to surpass the token-level encoder-only baseline (mBERT) on various downstream multilingual tasks. ByT5~\citep{xue2022byt5} investigated encoder-decoder models that operate directly at the byte level, opting not to use patching techniques. While these models demonstrated greater robustness to noise and achieved competitive results with four times less training data than tokenizer-based approaches, the absence of patching compelled the model to perform attention operations across every individual byte, incurring severe computational costs. Using bytes as input instead of subword tokens inherently lengthens input sequences, creating scalability challenges for byte-level models. More recently, \cite{wang2024mambabyte} leveraged the Mamba Architecture~\citep{gu2023mamba}, which retains a fixed-size memory state even for extremely long contexts. Without employing patching, they trained a byte-level Mamba model and, under controlled FLOP conditions at the 350M parameter scale, obtained superior bits-per-byte performance compared to byte-level transformer models across several benchmarks.

\textbf{Patching-based approaches:} Employing patching techniques can significantly reduce the excessive number of flops typically incurred by byte-level LLMs while still preserving model quality. Several studies have achieved promising outcomes in settings with modest model sizes and limited training data. For instance, \cite{nawrot-etal-2022-hierarchical} introduced the hourglass transformer, which utilizes static patching methods for both downsampling and upsampling; at 150M parameters, this architecture surpasses other byte-level baselines. Building upon this, \cite{nawrot-etal-2023-efficient} advance the methodology by implementing dynamic patching strategies. These include a boundary-predictor trained via end-to-end methods, another boundary-predictor with supervision informed by tokenizers, and an entropy-based approach similar to {\textsc BLT}. Their results indicate that such techniques can yield better performance than contemporary vanilla transformer models on language modeling benchmarks, specifically for 40M parameter models trained on roughly 400M tokens. In addition, \cite{lester2024training} explore the impact of compressing sequences with arithmetic coding, achieving compression efficiencies exceeding those possible with BPE. Utilizing an equal-info windows approach, their model shows superior results over byte-level models in language modeling tasks, although it still lags behind systems based on subword representations.

Our research builds upon and is most directly informed by MegaByte~\citep{yu2023megabyte}, a decoder-only causal LLM architecture that utilizes a predetermined scheme for patching and concatenates byte representations into patches. This approach employs a local model in the decoder to reconstruct bytes from these patches. The creators of MegaByte illustrate that their model performs comparably to tokenizer-based systems with 1B parameters on a corpus of approximately 400B bytes. Across all our experiments, we include ablations of MegaByte and observe that static patching consistently underperforms relative to the most advanced compute-efficient tokenizer models within a fixed FLOP budget; our findings show that {\textsc BLT} successfully addresses this discrepancy. \cite{slagle2024spacebyte} independently report similar findings regarding MegaByte, proposing adaptations that include patching based on whitespaces and other whitespace-like characters, as well as the addition of a localized encoder. Their results indicate gains over token-based Transformer models within a constrained computation environment for certain domains, notably Github and arXiv, at the 1B parameter level. We extend our empirical study to encompass this variant as well, revealing that additional innovations in architecture are necessary to further scale byte-level models and fully reach parity with leading tokenizer models like Llama 3.

\section{Limitations and Future Work}

For the purposes of making architectural decisions in this study, we train our models using the optimal number of steps identified for Llama~3~\citep{dubey2024llama}. It should be noted, however, that these scaling laws were originally established for BPE-level transformers and might not produce ideal (data, parameter size) ratios when applied to {\textsc BLT}. Calculating scaling laws specifically tailored to {\textsc BLT} remains an open direction for future research and could potentially yield even more favorable scaling dynamics for our model. Furthermore, most of our experiments were limited to model sizes of up to 1B parameters, and it remains to be seen whether the best architectural configurations will shift as we extend our investigations to models with 8B parameters or larger, possibly revealing enhanced performance at such scales.

Current transformer libraries and code repositories are optimized predominantly for the high efficiency of tokenizer-centric transformer models. Although we provide theoretically matched FLOP experiments and utilize advanced implementations like FlexAttention to accommodate non-standard transformer layers, the actual wall-clock performance of our implementations may still lag behind that of tokenizer-based models and could be improved through additional optimization efforts.

Whereas {\textsc BLT} employs an independently trained entropy model for patching, an intriguing avenue for further research would involve developing the patching model through end-to-end training. Section \ref{sec:distillation} provides preliminary results suggesting the feasibility of converting tokenizer-based architectures—such as Llama 3, which is trained on over 10T tokens—into byte-level models. This is accomplished by initializing and keeping the global transformer's weights fixed. Continued investigation in this area could reveal approaches that preserve the advantages inherent to bytefying, and potentially achieve superior performance compared to the original tokenizer-based models, without necessitating full retraining from scratch.

\section{Conclusion}
\label{section:conclusion}

In this work, we introduce the Byte Latent Transformer (\textbf{BLT}), a novel model architecture that challenges the traditional reliance on static, fixed-vocabulary tokenization methods used in most large language models. {\textsc BLT} adopts a trainable, flexible approach for assembling bytes into variable-length patches, thereby enabling a data-dependent allocation of computational effort according to input complexity. This mechanism yields notable gains in efficiency and resilience. Through a thorough scaling analysis, we show that {\textsc BLT} achieves performance on par with tokenization-based architectures such as Llama 3, even at large scales (up to 8B parameters and 4T bytes processed), while delivering as much as a 50\% reduction in inference computational cost for minor decreases in standard metrics. Additionally, {\textsc BLT} introduces a novel scalability axis, wherein both model size and patch size can be grown concurrently within the same inference resource constraints—an advantage particularly relevant for computation budgets typical of practical applications. By operating directly at the byte level, {\textsc BLT} also extends the model’s capacity to process rare or noisy data and develops richer representations of sub-word features, further strengthening robustness. Collectively, these findings establish {\textsc BLT} as a compelling and scalable alternative to conventional tokenization, supporting more flexible, efficient, and robust language modeling.

\end{document}